\paragraph{Training and choosing an optical representation.}
Each fit started from the same visual metric and learned prototypes for the same regular-print inventory. Independent random-number generators preserved the sequence of identities, fonts and translations while changing the candidate images. For three quarters of views in expectation we rescaled both axes independently in $[0.78,1.22]$, interpolated ink and binarized it at one of three thresholds (0.32, 0.50 or 0.68). If all ink disappeared the original translated glyph was retained. We counted conflicting labels on the actual packed training images. For labelled examples $x$ with rendered observations $g(x)$ and targets $t(x)$, an exact reader on the observations exists if and only if $g(x)=g(y)$ implies $t(x)=t(y)$. Sufficiency follows by choosing the target of a representative of each nonempty fibre. Lean verified both implications. Existence alone provides neither a sample-complexity bound nor a procedure for generalizing to a new observation.

A subsequent reader chooses its augmented correspondence whenever fewer than four fifths of nonblank occurrences in either the source or the request exactly match its fixed image inventory. All-white strips contribute no evidence of novelty. Generation always retains the original source addresses. We assessed both optical routes on paired Chinese reading questions, historical-image requests and changes in the source evidence. The translation material came from XQuAD (Artetxe, Ruder and Yogatama, 2020), a professionally translated subset of SQuAD development. We paired questions with disjoint recorded answers in the visible Chinese passage and rendered each question in Chinese and English while keeping the Chinese answer references. Font coverage was required for both requests. Earlier local English SQuAD use was audited by question identity and paragraph digest. The metric requires a stopped answer image matching a whole reference, permitting its glyphs to appear in canonical or incoming print.
